% Source: https://artofproblemsolving.com/community/c3h1498075_simplifying_rational
% Stable identity: aops-topic-1498075-q01
\textbf{(a)} Take $a=b=1$: the left side is $\dfrac{1}{2}$ and the right side is $1-1=0$. So the expressions are not identically equal. (No algebraic trick, such as multiplying by a conjugate, can turn one into the other.)

\textbf{(b)} Put $t=\left(\dfrac {a}{b}\right)^2>0$ and divide the numerator and the denominator by $b^2$:
\[
\dfrac{b^2}{a^2+b^2}=\dfrac{1}{t+1},\qquad \dfrac{b^2}{a^2}-1=\dfrac{1}{t}-1=\dfrac{1-t}{t}.
\]
Equality means $t=(1-t)(1+t)=1-t^2$, i.e.\ $t^2+t-1=0$. The positive root is $t=\dfrac{\sqrt5-1}{2}$, so
\[
\boxed{\dfrac {a}{b}=\sqrt{\dfrac{\sqrt5-1}{2}}}\approx0.786 .
\]

\textbf{(c)} For example,
\[
\dfrac{b^2}{a^2+b^2}=\dfrac{1}{1+\left(\dfrac {a}{b}\right)^2}=1-\dfrac{a^2}{a^2+b^2}.
\]
